% Source: https://artofproblemsolving.com/community/c3h1869323_supertasking
% Stable identity: aops-topic-1869323-q01
Ball $k$ is put in at step $\left\lceil\dfrac{k}{10}\right\rceil$ and is taken out at step $k$. Since $\left\lceil\dfrac{k}{10}\right\rceil\le k$, ball $k$ is always in the urn when we want to take it out, and it is never put in again.

(a) After step $n$ we have put in $10n$ balls and taken out $n$ of them, so the urn contains $\boxed{9n}$ balls.

(b) Every ball $k$ is taken out at step $k$ and never returns, so no ball satisfies the condition. The urn is empty at noon: there are $\boxed{0}$ balls.

This is the Ross--Littlewood paradox: the counts $9,18,27,\dots$ grow without bound, yet the urn is empty at noon. There is no contradiction, because the number of balls at noon is not the limit of the numbers $9n$. The answer depends on \emph{which} balls are taken out: if at step $n$ we took out ball $10n$ instead, the balls $1,2,\dots,9,11,\dots$ would stay forever and infinitely many balls would remain.
